\paragraph{Conjugate-coordinate class times a fuzzy factor: closed forms of the level-one forms, Lichnerowicz reduction of the mixed term, the range of the background torsion and the coefficient on the null cone (conditional; bounded no-go; refinement of three statements of the preceding paragraph).}
This paragraph supplements the preceding one, in which the forms were sampled numerically at points, the range \(M\sim\coth^2\tau\) was not computed and isolated data of \(K\) were not excluded. Keep its setting and notation and write \(\vartheta=\coth^2\tau\).

\emph{Closed forms.} At \(y^i=0\) the Gram matrix of the \(s_i\) is \(G_{ij}=a\,\delta_{ij}-(a-b)\hat q_i\hat q_j\), with \(a=\ell^2q^2+(q\cdot y)^2\), \(b=-m^2q^2\) and \(\hat q\) the unit spatial direction of \(q\). Put \(x=a/b\): then \(x\le-\ell^2/m^2\) for space-like \(q\), \(x\ge1\) for time-like \(q\), \(|x|\to\infty\) on the null cone at fixed \(q\cdot y\neq0\), and \(y_\perp=y-(q\cdot y)q/q^2\) has \(y_\perp^2=m^2(x-1)\). With \(\bar h^{\mu\nu}=\tfrac{2m}3D_{ij}s_i^\mu s_j^\nu\) the fibre-averaged level-one torsion is \(-\tfrac12(q^a\bar h^{b\mu}-q^b\bar h^{a\mu})\) and \(\mathrm{EH}(\bar h)=-q^2\bar h\cdot\bar h/8m^6\). On the five traceless polarisations count the helicity about \(y_\perp\) in the three-space orthogonal to \(q\), and put \(\varphi=(\bar hy)^2/(\bar h\cdot\bar h\,y_\perp^2)\), so that \(\varphi=0,\tfrac12,\tfrac23\) for helicity \(\pm2,\pm1,0\). The fourth moment of the fibre vector gives exactly, for the eigenvalues relative to \(\mathrm{EH}(\bar h)\),
\[
\mathrm{IX}_2:\ -\tfrac65\bigl[2+(2+\tfrac1x)(1+(x-1)\varphi)\bigr],\qquad
\mathrm{SC}_2:\ \tfrac1{20}\bigl[4+2r+(2+\tfrac1x)(1+(x-1)\varphi)\bigr],
\]
with \(r=\tfrac12(x+\tfrac1x)\) on helicity \(\pm1\) and \(r=1\) otherwise, and the Gaussian moments give, as rational identities in \(q,y,m\),
\[
Q_2^D/m^4:\ \tfrac32\vartheta+(x-1)\varphi,\qquad
(Q_2^V-2Q_2^D)/m^4:\ 9\vartheta+(x-1)(8-6\varphi).
\]
The exact forms agree with these at eighteen sampled points. Hence the further form of the preceding paragraph is \(q^2(\bar hy)\cdot(\bar hy)\), the Maxwell form of \(a_b=\bar h_{b\nu}y^\nu\), that is the square of the contraction of the level-one torsion with \(y\); the double eigenvalue \(\tfrac32\vartheta\) belongs to helicity \(\pm2\); and \(\lambda=-\tfrac65(2+\tfrac1x)\), a function of \(x\) that changes sign at \(x=-\tfrac12\), reached only where \(m^2\ge2\ell^2\). The term \(8(x-1)\) is not a multiple of the Einstein--Hilbert form, since \((x-1)q^2\bar h\cdot\bar h=[y^2q^2-(q\cdot y)^2]\bar h\cdot\bar h/m^2\).

\emph{Dependence on the wave vector.} A coefficient on the five for every \(q\) at a base point must not depend on the helicity or on \(x\). At first order in the torsion on fibre-constant modes the \(K\)-weights of \(Q_2^D\) and \(Q_2^V-2Q_2^D\) do not depend on \(x\). Then helicity \(\pm1\), through the term \((x-1)^2/20x\) of \(\mathrm{SC}_2\), requires \(W_R^{(1)}=0\), and helicities \(0\) and \(\pm2\), through \((x-1)(2+\tfrac1x)\), then require \(W_X^{(2)}=0\). Since \(W_X^{(2)}=-\tfrac18\int\mathrm ds\,s^{-2}\tau^KB_0\) with \(\tau^K\ge0\) and \(B_0>0\), \(B_0\) being the sum \(B_i\) of the preceding paragraph with the weight \(\int_0^1e^{-s\Delta a}\,\mathrm da\) of the torsion vertex, this fails for every non-zero flux of \(K\). At this order and at the frozen symbol the isolated data of \(K\) left open in the preceding paragraph are excluded, and the value \(M^*\) quoted there belongs to one class of wave vectors: at \((5,3)\), \(n=30\), \(L/c=0\) in the single-scale model the zero of the first-order coefficient is \(M^*\simeq0.292\) at \(x\simeq-2.70\) and \(M^*\simeq0.223\) at \(x\simeq1.33\).

\emph{Lichnerowicz reduction of the mixed term.} For fibre-constant loop modes with projected coefficients (the transitions to \(J'=2\) enter through the weights \(B_i\) of the preceding paragraph, not here) and at the order considered, the operators are built from the averaged frame fields \(\bar E^a\), which are divergence free for \(\rho=1/m\); the second-harmonic product \(\langle\delta E\delta E\rangle\) carries fewer powers of \(q\). Then \(T^{ab\mu}=[\bar E^a,\bar E^b]^\mu\) and \(\partial_\mu(\rho T^{ab\mu})=0\) identically, so with \(g^{\mu\nu}=\eta_{ab}\bar E^{a\mu}\bar E^{b\nu}\), \(\phi=\tfrac12\ln(\rho/\sqrt{|g|})\) and \(B^\mu=-\Sigma_{ab}T^{ab\mu}\) every species has connection \(\partial\phi-\tfrac12B\) and endomorphism \(E_0-\tfrac14B_\mu B^\mu\), \(E_0=-\Box_g\phi-|\partial\phi|^2\), without a term linear in \(T\). The Dirac operator \(\Gamma_a\bar E^a\) is similar through \((\rho/\sqrt{|g|})^{1/2}\) to the Levi-Civita Dirac operator of \(g\) up to the axial torsion, and for any such frame, with \(Z(k)=Z^\mu k_\mu\), the index density satisfies
\[
\tfrac14g_{\mu\nu}\operatorname{tr}_DZ^\mu Z^\nu=-R[g]+4\Box_g\phi+4|\partial\phi|^2\mp\tfrac18A_aA^a,
\]
with \(A_d=\epsilon_{abcd}T^{abc}\) the axial torsion and the upper sign for \(\eta\), the lower for \(\delta\) (exact rational arithmetic, ten random two-jets each). On the background \(T^{ab\mu}=\bar T^{ab\mu}\) is the Levi-Civita spin connection of \(G=\eta/m^2\), \(R[G]=18\vartheta-24\) and \(R/6+E_0=2-\tfrac34\vartheta\) (with \(\vartheta=-y\cdot y/m^2\) off \(y^i=0\)), and \(\bar{\mathcal Z}_{ab}\) is \(\pm J_{ab}\), which commutes with \(\Box\), so that, without mass terms, \(e^{-sH_i}=e^{s\Box}e^{\pm sJ_{ab}\Sigma_i^{ab}}\) and \(H_i=C_{31}^{\mathrm{tot},i}-C_{41}-C_{\Sigma_i}\) for the vector and for the Dirac index (polynomial Poisson algebra, both signs of the bracket). The background torsion is therefore not an independent expansion field of the mixed term. Where the axial torsion and \(\delta\phi\) vanish at first order, as on the five for the fibre average at the leading symbol, the mixed term is a functional of \(g\) and \(\phi\) alone. Its order \(s\) is \(-\sqrt{|g|}R[g]\) at order \(q^2\), that is \(\tfrac12\mathrm{IX}_1=-\mathrm{EH}(\bar h)\) for a wave averaged over its phase, in agreement with the density \(\tfrac12\rho^2X\) of a Dirac field; the forms \(Q_i\) and the Dirac trace of the preceding paragraph and of this one are taken at unit amplitude of \(\delta T_0\) and are twice their phase averages, so that \(Q_1=m^4\mathrm{IX}_1\). Its order \(s^2\), up to total derivatives, is \(\tfrac1{72}R^2-\tfrac1{24}R_{\mu\nu\rho\sigma}R^{\mu\nu\rho\sigma}-2(\tfrac16R+E_0)^2\), a density quadratic in the curvature and \(E_0\), whose second variation meets the Einstein--Hilbert form only through background values of order \(\vartheta\); the form \(Q_2^D\) is only its part of fourth order in \(Z\) (at the background \(\tfrac{3125}{648}\) against the complete \(\tfrac{1783}{648}\) at \((5,3)\)), so the first background-torsion correction of the preceding paragraph is that of the frozen symbol and not the complete term of order \(s^2\). Thus the equality of the Dirac cone value with the coefficient of the Einstein term in the index density, recorded above as an equality of values and not derived there, is the Lichnerowicz identity with the antiperiodic gluing of Section~\ref{sec:status-B}, as for the Levi-Civita Dirac field of~\cite{Kabat1995}. The cross term of the same density with the conical deltas of Section~\ref{sec:status-B}, with \(\rho\) regular, gives the surface density \((1-n)\pi[\tfrac19R-\tfrac13R_{ijij}-\tfrac83(\tfrac16R+E_0)]\), which matches the Dirac weights stated there in the relative normalisation of \(\sigma_{\log}\) (the spin term there has a heat-trace prefactor \(4\pi\) times that of the terms in \(c_B\)); to this order the cone of the mixed term is not an independent quantity. The vector index does not reduce in this way, but it occurs only in the first term, whose weighted coefficient of order \(s^2\) is the analogue of \(\tfrac12\mathcal S_4\) of Section~\ref{sec:status-B}, with no term linear or quadratic in the curvature, so that the first-term form of the preceding paragraph is complete at that order. At a fibre point the frame has axial torsion and \(\hat X=m^{-4}q^2(v\cdot v)(\alpha\cdot\alpha)=-4\,\mathrm{EH}[h_n]+m^{-4}q^2(\alpha\cdot v)^2\) for \(h_n=m(\alpha\otimes v+v\otimes\alpha)\) (exact at six fibre points), half of the coefficient \(-4\) being axial. The identity above has the form of Equation~(3.31) of~\cite{MantaSteinackerTran2024} (the normalisation of the dual is not compared). As a comparison of structure only, the fibre trace of the preceding paragraphs corresponds to the higher-spin trace of the torsion term there, in the regime assumed there of a scale of \(K\) above the higher-spin masses, which that work replaces by one positive scale. The sections compared make no statement of one coefficient on the five polarisations: that is a criterion of this ledger, which is neither established in that work nor contradicted by it.

\emph{The range \(M\sim\vartheta\).} The frozen integral can be done: with \(\xi=s\vartheta/2\), \(\langle\operatorname{tr}_De^{-sZ}\rangle=4(1-\xi)e^{-\xi/2}\), \(\langle\operatorname{tr}_Ve^{-sZ}\rangle=2+2(1-4\xi)e^{-2\xi}\), and the second variation of the Dirac trace is \(-2s\,e^{-\xi/2}[1-(1+\nu\varphi)\xi+(\tfrac2{15}+\tfrac15\nu\varphi)\xi^2]\) times \(m^4\mathrm{EH}(\bar h)\), with \(\nu=(x-1)/\vartheta\) (exact; a direct evaluation agrees). For \(\tau^K\propto s^2e^{-sM}\) the weight \(M^{-2}\) of the preceding paragraph is replaced by \(F(M)=\tilde M^{-2}-\tfrac15\vartheta\tilde M^{-3}>0\), \(\tilde M=M+\tfrac14\vartheta\), which is bounded at \(M=0\); at \((5,3)\) the zeros at \(L/c=0\) are then absent for the sampled \(2\le n\le300\) and values of \(x\), and those at \(L/c=1,10\) for \(n=30,300\) and two values of \(x\) move by less than \(2\%\). This is a statement about the frozen integral only. The Levi-Civita form above shows that the terms dropped by freezing are of the same order in \(s\vartheta\): the curvature of \(G\) is of order \(\vartheta\). A frozen plane wave of two-derivative order needs the squared components of \(mq\) in the rest frame of \(y\) above that curvature and below the scale of \(K\), which implies \(M\gg\vartheta\) (a scale argument; the invariant \(m^2q^2\) vanishes on the null cone and is not the controlling quantity). In the range \(M\lesssim\vartheta\) no local two-derivative coefficient is defined by these forms, and the averaged-frame reading of the preceding paragraph lies there.

\emph{Null cone.} As \(q^2\to0\) at fixed \(q\cdot y\neq0\), \(\hat y_\perp\) tends to the direction of \(q\), helicity \(\pm2\) are the two transverse traceless polarisations with \(\bar h_{\mu\nu}y^\nu=0\), and their eigenvalues stay finite and equal: \(\mathrm{IX}_2\to-\tfrac{24}5\) and \(\mathrm{SC}_2\to\tfrac25\) (the exact forms approach these at nearly null \(q\)), so at the leading symbol the pole stays on the null cone of \(G\). At zeroth order in the background torsion, which by the preceding part requires \(M\gg\vartheta\), and for every \(K\) and every \(n\), these two polarisations therefore have one coefficient, the same for every null \(q\) (at fixed \(q\) the equality of the two values is the rotation about \(\hat q\))
\[
c_T^{\mathrm{TT}}=\tfrac16W_R^{(0)}+\tfrac7{30}W_R^{(1)}-2W_X^{(0)}-\tfrac{14}5W_X^{(2)},
\]
outside the hypothesis of five polarisations of the preceding paragraph; the other three keep distinct \(x\)-dependent values, so the tensor equation of the first law is not obtained. Eliminating the inert amplitudes through the induced forms alone, without the tree-level form, changes only helicity \(\pm1\) at three sampled cases, so it leaves this coefficient unchanged there. The first term shifts the cone: the helicity-two value \(8(x-1)\) contributes \((q\cdot y)^2\bar h\cdot\bar h\), and in the single-scale model the shift is of relative order \((L/c)\vartheta/(M^2W)\) with \(W=\sum_J(2J+1)/(M+\mu_J)\).

\emph{Sign and replicas.} The generators are traceless, so \(v^K\ge0\) as well as \(\tau^K\ge0\) by the Peierls--Bogoliubov inequality for Hermitian couplings (random non-commuting data), and \(B_0\ge B_1>0\) pointwise, \(B_1\) being the sum \(B_i\) of the preceding paragraph with the weight \(6\int_0^1a(1-a)e^{-s\Delta a}\,\mathrm da\) of the term \(F_1\) there. With \(V_i,T_i\) the integrals of \(s^{-2}v^K\), \(s^{-2}\tau^K\) against \(D\) for \(i=0\), \(B_1\) for \(i=1\), \(B_0\) for \(i=2\) (\(V_i,T_i\ge0\), assumed finite, and \(T_2\ge T_1\)),
\[
\tfrac16W_R-2W_X=\tfrac16V_0+\tfrac1{12}T_0>0,\qquad
c_T^{\mathrm{TT}}=\tfrac16V_0+\tfrac1{12}T_0+\tfrac7{30}(V_1-T_1)+\tfrac7{20}T_2>0
\]
for a non-zero flux: the sign of \(W_R\) is not fixed, but that of these coefficients at that order is, and no sign is claimed for \(M\lesssim\vartheta\). For commuting fluxes with \(f_1=f_2+f_3\) on every mode, \(v^K=\tau^K\), so that all \(W_R\) vanish. Against \(c_T^{\mathrm{TT}}\) the periodic and smooth replicas give \(c_A-c_T^{\mathrm{TT}}<0\) for a non-zero flux: \(-\tfrac14T_0-\tfrac7{30}(V_1-T_1)-\tfrac7{20}T_2\) and \(-\tfrac14V_0-\tfrac7{30}(V_1-T_1)-\tfrac7{20}T_2\). The mixed part of the fixed-state gap is \(-\tfrac7{30}(V_1-T_1)-\tfrac7{20}T_2<0\), so equality on these two polarisations needs the cutoff-dependent first-term cone contribution of the fixed-state replica, or contact and edge terms, to supply it.

These statements are for the hypotheses of the preceding paragraph, with the fibre-constant reduction for the Lichnerowicz form and first order in the torsion on fibre-constant modes for the exclusion of isolated data; the cross term is dropped, which on the null cone requires the squared components of \(mq\) to lie below the scale of \(K\). They do not compute the second variation of the order \(s^2\) density on the leaf, the torsion corrections on modes with \(J\ge1\), the fibre-flux part, the elimination of the inert amplitudes with the tree-level form, the trace polarisation or a Newtonian limit; they do not construct \(K\), do not make the Lorentzian evaluation more than formal, and do not identify \(\bar h\) with the emergent metric of Chapter~\ref{ch:vierbein-emergence} or bear on Assumption~\ref{ass:lorentzian}, so no value of \(G_{\mathrm{eff}}\) and no sign of it for the emergent metric follows. Problems~B and~F remain open and PROGRAMME-CONDITIONAL.
